\def\stat{pavelieva}








\def\tit{ПОИСК СООТВЕТСТВИЙ МЕЖДУ КЛЮЧЕВЫМИ ТОЧКАМИ ИЗОБРАЖЕНИЙ РАДУЖНЫХ ОБОЛОЧЕК ГЛАЗ 


С~ПОМОЩЬЮ МЕТОДА ПРОЕКЦИОННОЙ ФАЗОВОЙ КОРРЕЛЯЦИИ$^*$}





\def\titkol{Поиск соответствий между ключевыми точками изображений радужных оболочек глаз} 


%с~помощью метода проекционной фазовой корреляции} 





\def\autkol{Е.\,А. Павельева}





\def\aut{Е.\,А. Павельева$^1$}





\titel{\tit}{\aut}{\autkol}{\titkol}





{\renewcommand{\thefootnote}{\fnsymbol{footnote}}


\footnotetext[1]


{Работа выполнена при поддержке гранта Президента РФ МД-72.2001.9 и грантов РФФИ 10-07-00433


и~13-07-00438.}}











\renewcommand{\thefootnote}{\arabic{footnote}}


\footnotetext[1]{Московский государственный университет им.\ М.\,В.~Ломоносова, 


факультет вычислительной математики и кибернетики, 


paveljeva@yandex.ru} 





      \thispagestyle{headings}


      


                  \label{st\stat}











\Abst{Работа основана на анализе и сопоставлении наиболее информативных точек текстуры 


радужной оболочки глаза~--- ключевых точек. Для нахождения ключевых точек используется 


преобразование Эрмита. Предлагается осуществлять поиск соответствий между 


ключевыми точками радужных оболочек с помощью метода проекционной фазовой 


корреляции локальных окрестностей ключевых точек. В~методе проекционной фазовой 


корреляции преобразование Фурье применяется к разложениям функций интенсивности 


изображений в ряд Фурье по функциям Эрмита. В~предложенном методе сопоставления 


ключевых точек не используется информация, связанная с их расположением. 


Метод позволяет добиться хороших результатов сравнения изображений и избежать ошибок, 


связанных с локальными сдвигами частей изображений и с попаданием век, ресниц и бликов 


на изображения.}





\KW{распознавание по радужной оболочке глаза; фазовая корреляция; функции Эрмита; 


ключевые точки; биометрия}








\DOI{10.14357\!/08696527130206}





\section{Введение}





Выделение на изображении наиболее информативных ключевых точек является 


распространенным подходом к решению многих биометрических задач. По 


выделенным ключевым точкам составляется вектор свойств и производится 


сравнение изображений. В~случае распознавания по радужной оболочке глаза 


этот подход встречается в работах~[1--4].





В~работе~\cite{1-pav} ключевые точки выделяются с помощью фильтра Габора. 


В~качестве вектора свойств радужной оболочки вычисляются расстояния между 


ключевыми точками. Однако в случае попадания на изображение радужной 


оболочки бликов, век, ресниц может отсутствовать большое количество 


ключевых точек радужной оболочки.





В работе~\cite{4-pav} ключевые точки выделяются с помощью сверток с 


выбранной функцией преобразования Эрмита. Сравнение ключевых точек 


происходит с помощью расстояния Хэмминга. Изображение разбивается на 


блоки, и значение блока равняется единице, если там присутствует хотя бы одна 


ключевая точка, и равняется нулю в противном случае. Мерой близости 


изображений является число совпадающих блоков. Таким образом, 


используются соответствия между ключевыми точками, связанные лишь с их 


расположением. 





В данной работе предложено сопоставление ключевых точек радужных 


оболочек глаз методом проекционной фазовой корреляции локальных 


окрестностей ключевых точек. 





Фаза преобразования Фурье изображения несет в себе гораздо больше информации, чем его 


амплитуда~\cite{5-pav}. Метод фазовой корреляции~\cite{6-pav, 7-pav} использует 


только информацию о фазе дискретного преобразования Фурье для определения 


близости и сдвигов изображений. Дискретное преобразование Фурье 


подразумевает периодическое продолжение изображений, поэтому в случае 


изображений небольших размеров в методе фазовой корреляции могут 


возникать ошибки вычисления, связанные с ложным оконтуриванием. В~работе 


предлагается метод проекционной фазовой корреляции, в котором 


преобразование Фурье применяется к разложениям функций интенсивности 


изображений в ряд по функциям Эрмита. Метод проекционной фазовой 


корреляции не требует периодического продолжения изоб\-ра\-же\-ний для 


вычисления преобразования Фурье, поэтому имеет преимущество на 


изображениях небольших размеров.





\begin{figure}[b] %fig1


          \vspace*{-4pt}


 \begin{center}


 \mbox{%


 \epsfxsize=127.656mm


 \epsfbox{pav-1.eps}


 }


 \end{center}


 \vspace*{-11pt}


\Caption{Выделение радужной оболочки и маски радужной оболочки~(\textit{а}); 


нормализованное изображение радужной оболочки~(\textit{б}); нормализованное 


изображение после предобработки~(\textit{в})}


\end{figure}





\section{Предобработка изображений радужных оболочек глаз}





На изображении выделяется радужная оболочка глаза и маска радужной 


оболочки~--- области радужной оболочки, не закрытые ресницами, веками и 


бликами. Далее радужная оболочка переводится в прямоугольное 


нормализованное изображение размером $512\times64$~пикселов~\cite{4-pav}, 


происходит выравнивание освещенности и повышение контрастности 


изображения (рис.~1)~\cite{8-pav}.











\section{Метод ключевых точек}





После предобработки изображения радужной оболочки на нем выделяются 


ключевые точки при помощи свертки с выбранной функцией преобразования 


Эрмита.





Функции преобразования Эрмита~\cite{9-pav} задаются формулой:


$$


\varphi_n(x)=\fr{1}{\sqrt{2^n  n!}}\,\fr{1}{\sigma\sqrt{\pi}}\, 


e^{-(x/\sigma)^2} H_n\left( \fr{x}{\sigma}\right)\,,\enskip n=0,1,2,\ldots , 


$$


где $\sigma$~--- параметр масштаба; $H_n(x)$~--- полиномы Эрмита:


\begin{gather*}


H_0(x)=1\,;\qquad H_1(x)=2 x\,;\\


H_n(x)=2 x  H_{n-1}(x)-2 (n-1) H_{n-2}(x)\,.


\end{gather*}





\begin{floatingfigure}{65mm} %fig2


          \vspace*{1pt}


 \begin{center}


 \mbox{%


 \epsfxsize=50mm


 \epsfbox{pav-2.eps}


 }


 \end{center}


 \vspace*{-6pt}


\Caption{Функция преобразования Эрмита~$\varphi_{2,0}$}


\end{floatingfigure}





Двумерные функции преобразования Эрмита можно представить в виде 


произведения одномерных функций:


$$


\varphi_{m,n}(x,y)=\varphi_m(x) \varphi_n(y)\,.


$$











Текстура радужной оболочки глаза сильнее изменяется в горизонтальном 


направлении нормализованного изображения. Сверткой (или корреляцией) с 


функцией $\varphi_{2,0}$ (рис.~2) можно выделить текстурные особенности 


радужной оболочки глаза~\cite{8-pav}.











В каждой точке $(x_0,y_0)$ нормализованного изображения радужной оболочки 


вычисляется свертка данного изображения с функцией преобразования Эрмита 


$\varphi_{2,0}$ со значением параметра масштаба $\sigma\hm=3$: $(I(x,y)* 


\varphi_{2,0}(x,y))\vert_{(x_0,y_0)}$.





Далее выбираются точки со значениями сверток больше среднего значения 


модулей сверток и отстоящие друг от друга не менее чем на 3~пиксела (рис.~3).





Ключевыми точками считаются $N$ выбранных точек ($N\hm\leq 150$) с 


наибольшими значениями сверток.





\begin{figure} %fig3


          \vspace*{1pt}


 \begin{center}


 \mbox{%


 \epsfxsize=128mm


 \epsfbox{pav-3.eps}


 }


 \end{center}


 \vspace*{-6pt}


\Caption{Ключевые точки на исходном нормализованном изображении, полученные с 


помощью сверток с функцией преобразования Эрмита $\varphi_{2,0}(x,y)$, $\sigma\hm=3$}


\end{figure}





\section{Метод фазовой корреляции}





Метод фазовой корреляции~\cite{6-pav, 7-pav} позволяет вычислять меру 


близости изоб\-ра\-же\-ний. Для двух дискретных изображений $f(x,y)$ и $g(x,y)$ 


размером $M\times N$ вычисляются их дискретные преобразования Фурье 


$F(u,v)$ и $G(u,v)$, и для спектральных функций $F(u,v)$ и $G(u,v)$ 


вычисляется их взаимный фазовый спектр:


$$


R_{FG}(u,v)= \fr{F(u,v) \overline{G(u,v)}}{ \vert F(u,v) 


\overline{G(u,v)}\vert} = e^{i(\varphi_{F}(u,v)-\varphi_{G}(u,v))}\,,


$$


т.\,е.\ спектральная функция с единичным модулем, фаза которой равна 


разности фаз функций~$F$ и~$G$. Далее от функции взаимного фазового 


спектра вычисляется обратное дискретное преобразование Фурье и получается 


функция фазовой корреляции (POC (phase-only correlation) функ\-ция):


$$


\mathrm{POC}_{fg}(x,y) =F^{-1}\left[ R_{FG}(u,v)\right]\,.


$$





В случае <<похожих>> изображений POC-функция имеет один четкий пик, 


высота пика определяет меру близости изображений, а положение пика 


соответствует смещению одного изображения относительно другого. Если же 


функции <<не похожи>>, то POC-функ\-ция не дает четкого пика. 


В~работе~\cite{10-pav} показано, что если брать обратное преобразование 


Фурье не от всего спектрального сигнала, а только от его части, 


соответствующей низким частотам, то пик POC-функ\-ции получается более 


четким и устойчивым.





\section{Метод проекционной фазовой корреляции}





В методе проекционной фазовой корреляции~\cite{10-pav} преобразование Фурье 


применяется к разложениям исходных непрерывных функций интенсивности 


изображений в ряд по функциям Эрмита $\psi_{m,n}(x,y)$~\cite{11-pav}.





Одномерные функции Эрмита задаются следующей формулой:


$$


\psi_n(x) =\fr{1}{\sqrt{2^n n!}}\, \fr{1}{\sqrt{\sigma\sqrt{\pi}}}\, 


e^{-x^2/(2\sigma^2)} H_n\left( \fr{x}{\sigma}\right)\,,\enskip n=0,1,2,\ldots


$$





Функции Эрмита являются собственными функциями преобразования Фурье 


$$


F[f]\hm=\fr{1}{\sqrt{2\pi}}\int\limits_{-\infty}^\infty f(x) e^{-i\lambda x}\,dx


$$ 


с собственными значениями $\pm 1, \pm i$: 


$$


F[\psi_n]\hm= (-i)^n\psi_n


$$


 и  образуют полную ортонормированную в пространстве $L_2 (R)$ систему функций.





Двумерные функции Эрмита можно представить в виде произведения 


одномерных:


$$


\psi_{m,n}(x,y)=\psi_m(x) \psi_n(y)\,.


$$





Пусть исследуемые функции $f(x,y)$ и $g(x,y)$ заданы в прямоугольнике $D=[-


A,A]\times [-B,B]$. Продолжим их нулем в $R^2\backslash D$ и рассмотрим 


частичную сумму разложения в ряд по функциям Эрмита $\psi_{l,j}(x,y)$, 


$l\hm= \overline{0,m}$, $j\hm=\overline{0,n}$:


\begin{equation}


f(x,y)\approx f_1(x,y) =\sum\limits_{l=0}^m \sum\limits_{j=0}^n c_{lj} 


\psi_{l,j}(x,y)\,,


\label{e1-pav}


\end{equation}


где коэффициенты Эрмита


$$


c_{lj} =\iint\limits_{R^2} f(x,y) \psi_{l,j}(x,y)\,dxdy\,.


$$





Тогда аппроксимация преобразования Фурье примет следующий вид:


$$


\mathrm{HF}[f,m,n] =F[f_1] =f\left[ \sum\limits_{l=0}^m \sum\limits_{j=0}^n c_{lj} 


\psi_{l,j}(x,y) \right] = \sum\limits_{l=0}^m \sum\limits_{j=0}^n c_{lj} (-i)^{l+j} 


\psi_{l,j}(x,y)\,.


$$


Аналогично вычисляется аппроксимация преобразования Фурье для $g(x,y)$. 


Далее для спектральных функций $\mathrm{HF}[f,m,n]$ и $\mathrm{HF}[g,m,n]$ находится их 


взаимный фазовый спектр $R_{fg}(x,y)$ и вычисляется функция проекционной 


фазовой корреляции $\mathrm{HPPOC}_{fg}(x,y)$. Поскольку $F^{-1}[R_{fg}(x,y)]\hm= 


F[R_{fg}(-x,-y)]$, то


$$


\mathrm{HPPOC}_{fg}(x,y)= \mathrm{HF}\left[R_{fg}(-x,-y),m,n\right]\,.


$$





\section{Сравнение метода фазовой корреляции и~метода 


проекционной фазовой корреляции}





Угол поворота одного глаза относительно другого соответствует циклическому 


сдвигу одного нормализованного изображения относительно другого. Будем 


искать сдвиг методами фазовой корреляции (POC) и проекционной фазовой 


корреляции (HPPOC). В~случае изображений одной радужной оболочки 


(рис.~4) пики POC- и HPPOC-функ\-ций выражены четко, при этом позиция 


пика соответствует сдвигу второго изображения относительно первого 


(повороту второго глаза относительно первого), а высота пика определяет меру 


близости функций. В~случае изображений разных радужных оболочек (рис.~5) 


нет четкого пика POC- и HPPOC-функ\-ций.





\begin{figure}[b] %fig4


          \vspace*{-6pt}


 \begin{center}


 \mbox{%


 \epsfxsize=128.161mm


 \epsfbox{pav-4.eps}


 }


 \end{center}


 \vspace*{-9pt}


\Caption{Нормализованные изображения глаз одного человека~(\textit{а}, \textit{б}), их 


POC-функ\-ция~(\textit{в}) и HPPOC-функ\-ция~(\textit{г}). Чем выше значения функций POC 


или HPPOC, тем темнее соответствующие области изображений для POC и HPPOC}


\end{figure}











Однако при небольших размерах изображений метод проекционной фазовой 


корреляции дает более точные результаты, чем метод фазовой корреляции 


(рис.~6). Это связано с тем, что метод HPPOC не требует периодического 


продолжения изображений для вычисления преобразования Фурье (в отличие от 


метода POC), и поэтому в методе HPPOC не возникает ошибок от скачков на 


границах (ложного оконтуривания). В~случае изображений маленького размера 


ошибки в методе POC, связанные с ложным оконтуриванием, могут приводить к 


неверным результатам.





\begin{figure} %fig5


          \vspace*{1pt}


 \begin{center}


 \mbox{%


 \epsfxsize=128mm


 \epsfbox{pav-5.eps}


 }


 \end{center}


 \vspace*{-9pt}


\Caption{Нормализованные изображения глаз разных людей~(\textit{а}, \textit{б}), их 


POC-функ\-ция~(\textit{в}) и HPPOC-функ\-ция~(\textit{г})}


\vspace*{6pt}


\end{figure}





\begin{figure} %fig6


          \vspace*{1pt}


 \begin{center}


 \mbox{%


 \epsfxsize=127.795mm


 \epsfbox{pav-6.eps}


 }


 \end{center}


 \vspace*{-9pt}


\Caption{Части изображений одного глаза (\textit{а}, \textit{б}); POC-функ\-ция~(\textit{в}); 


HPPOC-функ\-ция~(\textit{г})}


\end{figure}





\section{Сопоставление ключевых точек радужных оболочек глаз 


методом проекционной фазовой корреляции}











Предлагается метод сопоставления ключевых точек радужных оболочек глаз, 


основанный на анализе проекционной фазовой корреляции их 


окрестностей~\mbox{\cite{8-pav, 10-pav}}. В~каж\-дой ключевой точке выделяется 


локальная окрестность размера $25\times 10$~пикселов с центром в данной 


точке. Данное изображение раскладывается в сумму~(\ref{e1-pav}) по функциям 


Эрмита $\psi_{lj}$, $l\hm= 0,1,\ldots , 9$, $j\hm= 0,1,2,3,5$. Чем меньше функций 


Эрмита берется для параметризации, тем меньше коэффициентов будет 


храниться и участвовать в сравнении. Чем больше функций Эрмита, тем лучше 


будет аппроксимация функций.





\begin{figure} %fig7


          \vspace*{1pt}


 \begin{center}


 \mbox{%


 \epsfxsize=128mm


 \epsfbox{pav-7.eps}


 }


 \end{center}


 \vspace*{-6pt}


\Caption{Сопоставление ключевых точек двух изображений одного глаза методом 


проекционной фазовой корреляции: исходные нормализованные изображения (\textit{а}, 


\textit{б}, выделены две сравниваемые ключевые точки); окрестность (увеличенная) ключевой 


точки первого изображения~(\textit{в}); окрестность ключевой точки второго 


изображения~(\textit{г}); HPPOC-функция для данных окрестностей~(\textit{д})}


\end{figure}





\begin{figure} %fig8


          \vspace*{1pt}


 \begin{center}


 \mbox{%


 \epsfxsize=128mm


 \epsfbox{pav-8.eps}


 }


 \end{center}


 \vspace*{-6pt}


\Caption{Сопоставление ключевых точек (близких по месторасположению) изображений 


разных глаз методом проекционной фазовой корреляции}


\end{figure}











Сначала определяется угол поворота одного глаза относительно другого 


методом проекционной фазовой корреляции, примененным к целым 


нормализованным изображениям~\cite{10-pav} либо к правой верхней четверти 


нормализованных изображений для увеличения быстродействия работы метода. 


Далее каждая ключевая точка одного изображения сопоставляется с каждой 


ключевой точкой второго изображения из окрестности $40 \times 10$~пик\-се\-лов 


методом проекционной фазовой корреляции окрестностей (размером $25\times 


10$~пикселов) данных ключевых точек. Таким образом, считаем, что возможен 


поворот глаза на угол до $\pm 14^\circ$. Если отношение пика HPPOC-функции 


ко второму локальному максимуму этой функции больше значения $k\hm=3$, то 


считается, что ключевые точки соответствуют друг другу. Пример алгоритма 


сопоставления ключевых точек показан на рис.~7--9. 














В случае изображений одного глаза (см.\ рис.~7) для сопоставляемых ключевых 


точек есть четкий пик HPPOC-функ\-ции. В~случае изображений разных глаз 


(см.\ рис.~8) для сопоставляемых ключевых точек нет четкого пика 


HPPOC-функ\-ции.








На рис.~9 представлены нормализованные изображения одного глаза и показано 


их сопоставление методом проекционной фазовой корреляции. Черным цветом 


обозначены ключевые точки, не имеющие сопоставимых ключевых точек на 


другом изображении при использовании метода проекционной фазовой 


корреляции. Белые точки~--- точки, имеющие сопоставимые ключевые точки на 


другом изображении. Из рисунка видно, что чем меньше порог~$k$ отношения 


пика HPPOC-функ\-ции ко второму локальному максимуму этой функции, тем 


больше ключевых точек оказываются сопоставимыми.











\section{Сопоставление ключевых точек радужных оболочек глаз 


с~помощью~расстояния Хэмминга}











При сопоставлении ключевых точек с помощью расстояния 


Хэмминга~\cite{4-pav} изоб\-ра\-же\-ние разбивается на блоки $4\times 4$ пиксела, и 


значение блока приравнивается~1, если в него попадает хотя бы одна ключевая 


точка, и~0, если не попадает ни одной ключевой точки. Расстоянием между 


изображениями считается число совпадающих блоков. Таким образом, 


используются соответствия между ключевыми точками, связанные лишь с их 


расположением. В~случае локальных сдвигов изображений соответствующие 


друг другу ключевые точки изображений одного глаза могут попадать в 


соседние блоки. В~случае попадания на изображение радужной оболочки 


бликов, век, ресниц может отсутствовать большое количество ключевых точек, 


присутствующих на другом изображении той же радужной оболочки (рис.~10). 


Если же глаза разные, то в одной и той же области могут оказаться ключевые 


точки (т.\,е.\ соответствующие блоки будут совпадать), но соответствующие 


разным текстурным особенностям.





\begin{figure} %fig9


          \vspace*{1pt}


 \begin{center}


 \mbox{%


 \epsfxsize=128.149mm


 \epsfbox{pav-9.eps}


 }


 \end{center}


 \vspace*{-6pt}


\Caption{Сопоставление ключевых точек изображений одного глаза методом проекционной 


фазовой корреляции при различных порогах~$k$:


(\textit{а})~$k\hm=2$; (\textit{б})~3 (используется при анализе результатов работы метода в разд.~9);


(\textit{в})~$k\hm=4$}


\end{figure}





На рис.~10,\,\textit{а} и~10,\,\textit{б} приведены те же изображения, что и на рис.~9, и приведено 


сопоставление ключевых точек этих изображений с помощью расстояния 


Хэмминга (рис.~10,\,\textit{в}). Ключевые точки первого изображения обозначены черным цветом, 


второго~--- белым. Разобьем приведенные изображения на три части. Видно, 


что в правой части изображений ключевые точки на двух изображениях 


соответствуют друг другу. В~центре изображений заметно смещение ключевых 


точек изображений (это связано с локальным смещением центральных частей 


изображения). Слева на рис.~10,\,\textit{б} присутствует веко, поэтому там 


отсутствуют ключевые точки, присутствующие на рис.~10,\,\textit{а}.





\begin{figure} %fig10


          \vspace*{1pt}


 \begin{center}


 \mbox{%


 \epsfxsize=128mm


 \epsfbox{pav-10.eps}


 }


 \end{center}


 \vspace*{-6pt}


\Caption{Сопоставление ключевых точек изображений одного глаза методом Хэмминга}


\end{figure}














\section{Результаты алгоритмов сопоставления ключевых точек}





Для сравнения двух предложенных методов сопоставления ключевых точек 


используется часть базы данных CASIA-IrisV3~\cite{12-pav}, содержащая 


изображения глаз первых 20~людей в базе: 224~изображений, образующих 


40~множеств $M_1,\ldots , M_{40}$ (каждое множество содержит изображения 


одного из двух глаз одного из 20~людей). В~этой базе данных для каждой пары 


изображений сопоставляются ключевые точки с помощью расстояния Хэмминга 


и методом проекционной фазовой корреляции. Расстоянием между 


изображениями считается число совпадающих (по методу Хэмминга или методу 


проекционной фазовой корреляции) ключевых точек.





Возникает вопрос о выборе порогового значения: если расстояние между 


изоб\-ра\-же\-ни\-ями выше этого значения, глаза считаются одинаковыми, если 


ниже~--- разными. Ошибка I~рода, или FRR (False Rejection Rate),~--- значение в 


процентах, показывающее, в скольких случаях одинаковые глаза считаются 


разными. Ошибка II~рода, или FAR (False Acceptance Rate),~--- значение в 


процентах, показывающее, в скольких случаях разные глаза считаются 


одинаковыми. Величина EER (Equal Error Rate)~--- величина ошибки работы 


метода, при которой ошибки I и~II~рода равны.





При сопоставлении ключевых точек с помощью расстояния Хэмминга величина 


EER составила 3,43\%, при сопоставлении методом проекционной фазовой 


корреляции $\mathrm{EER}\hm=0{,}02\%$. На рис.~11 приведены результаты 


сопоставления ключевых точек с помощью расстояния Хэмминга~(\textit{а}) и с 


помощью проекционной фазовой корреляции~(\textit{б}). 


Свет\-ло-се\-рым цветом показано распределение всевозможных сравнений 


изображений одного глаза, тем\-но-се\-рым цветом~--- распределение 


всевозможных сравнений изображений разных глаз, черным~--- пересечение 


этих распределений.





\begin{figure} %fig11


          \vspace*{1pt}


 \begin{center}


 \mbox{%


 \epsfxsize=127.646mm


 \epsfbox{pav-11.eps}


 }


 \end{center}


 \vspace*{-6pt}


\Caption{Сопоставление ключевых точек с помощью расстояния Хэмминга~(\textit{а}) и 


проекционной фазовой корреляции~(\textit{б})}


\end{figure}





\section{Заключение}





Предложен метод сопоставления ключевых точек изображений радужных 


оболочек глаз на основе метода проекционной фазовой корреляции. Сравнение 


ключевых точек происходит по локальной информации окрестностей ключевых 


точек. В~методе не используются никакие соответствия между ключевыми 


точками, связанные с их взаиморасположением. Предложенный метод 


сопоставления ключевых точек достаточно устойчив к присутствию век, ресниц 


и бликов на изображении, а также к локальным сдвигам частей изображения.





{\small\frenchspacing


{%\baselineskip=10.8pt


\addcontentsline{toc}{section}{Литература}


\begin{thebibliography}{99}





\bibitem{3-pav} %1


\Au{Ma L., Tan T., Wang Y., Zhang D.} Efficient iris recognition by characterizing 


key local variations~/\!\!/ IEEE Trans. Image Process., 2004. Vol.~13. 


No.\,6. P.~739--750. 





\bibitem{1-pav} %2


\Au{Yu L., Zhang D., Wang K.} The relative distance of key point based iris 


recognition~/\!\!/ Pattern Recogn., 2007. Vol.~40. No.\,2. P.~423--430.


\bibitem{2-pav} %3


\Au{Hollingsworth K., Bowyer K., Flynn~P.} The best bits in an iris code~/\!\!/ IEEE 


TPAMI, 2009. Vol.~31. No.\,6. 


P.~964--973.





\bibitem{4-pav}


\Au{Павельева Е.\,А., Крылов А.\,С.} Поиск и анализ ключевых точек радужной 


оболочки глаза методом преобразования Эрмита~/\!\!/ Информатика и её 


применения, 2010. Т.~4. Вып.~1. С.~79--82.


\bibitem{5-pav}


\Au{Oppenheim A.\,V., Lim~J.\,S.} 


The importance of phase in signals~// IEEE Proceedings, 1981. Vol.~69. No.\,5. P.~529--541.


\bibitem{6-pav}


\Au{Miyazawa K., Ito K., Aoki T., Kobayashi~K., Nakajima~H.} A~phase-based iris 


recognition algorithm~/\!\!/ Advances in biometrics~/ 


Eds. D.~Zhang, A.\,K.~Jain. Lecture notes in computer science ser.~--- 


Berlin, Heidelberg: Springer, 2006.  Vol.~3832. 


P.~356--365.


\bibitem{7-pav}


\Au{Nagashima S., Ito K., Aoki T., Ishii~H., Kobayashi~K.} High-accuracy estimation 


of image rotation using 1D phase-only correlation~/\!\!/ IEICE Trans. Fund. Electr., 2009. 


Vol.~E92-A. No.\,1. P.~235--243.


\bibitem{8-pav}


\Au{Павельева Е.\,А.} Метод проекционной фазовой корреляции в ключевых 


точках радужной оболочки глаза~/\!\!/ ГрафиКон'2012: Труды 22-й Междунар. 


конф. по компьютерной графике и зрению.~--- М.: МАКС Пресс, 2012. 


С.~128--132.


\bibitem{9-pav}


\Au{Martens J.-B.} The Hermite transform-theory~/\!\!/ IEEE Trans. 


Acoust. Speech, 1990. Vol.~38. No.\,9. P.~1595--1606.


\bibitem{10-pav}


\Au{Павельева Е.\,А., Крылов А.\,С.} Определение локальных сдвигов 


изображений радужных оболочек глаз методом проекционной фазовой 


корреляции~/\!\!/ ГрафиКон'2011: Труды 21-й Междунар. конф. по 


компьютерной графике и зрению.~--- М.: МАКС Пресс, 2011. С.~188--191.


\bibitem{11-pav}


\Au{Krylov A., Korchagin D.} Fast Hermite projection method~/\!\!/ Image analysis and recognition~/ Eds.


A.~Campilho, M.\,S.~Kamel. Lecture notes in computer science ser.~---


Berlin, Heidelberg: Springer-Verlag, 2006. Vol.~4141. P.~329--338.


\bibitem{12-pav}


CASIA-IrisV3 database. Available at: {\sf http://www.cbsr.ia.ac.cn/IrisDatabase.htm}.


(accessed October 25, 2013).











 \end{thebibliography}


} }








\vspace*{12pt}





\hrule





\vspace*{2pt}





\hrule





\def\tit{THE SEARCH FOR MATCHES BETWEEN THE IRIS KEY POINTS USING HERMITE PROJECTION 


PHASE-ONLY CORRELATION METHOD}





\def\aut{E.\,A.~Pavelyeva}








\titelr{\tit}{\aut}





\vspace*{3pt}








\noindent


Faculty of Computational Mathematics and Cybernetics, M.\,V. Lomonosov Moscow 


State University,  Moscow, Russia, paveljeva@yandex.ru











\vspace*{3pt}





\Abste{The paper is based on the analysis and comparison of the iris key points which are the most 


informative points of the iris image texture. The key\linebreak\vspace*{-12pt}


}





\Abstend{points are selected using Hermite transform. The 


search for matches between the iris key points using Hermite projection phase-only correlation of 


local areas around key points is suggested. In Hermite projection phase-only correlation method, the 


Fourier transform is applied to the expansions of images intensity functions in a Fourier series of 


Hermite functions. The proposed method of key points comparisons does not use any information 


about key points locations, is robust to eyelids, eyelashes, and glares, and gives good results.}





\KWE{iris recognition; phase only correlation; Hermite functions; key points; biometrics}





\DOI{10.14357\!/08696527130206}





\vspace*{-9pt}





\Ack 





\vspace*{-3pt}





\noindent 


The research was supported by the President of Russian Federation (Grant No.\,MD-72.2001.9) and by the Russian


Foundation for Basic Research (Grants Nos.\,10-07-00433 and 13-07-00438).





\renewcommand{\bibname}{\protect\rmfamily References}





\vspace*{-3pt}





 {\small\frenchspacing


{%\baselineskip=10.8pt


%\addcontentsline{toc}{section}{Литература}


\begin{thebibliography}{99}





\vspace*{-3pt}





\bibitem{3-p-1} %1


\Aue{Ma, L., T.~Tan, Y.~Wang, and D.~Zhang}. 2004. Efficient iris recognition by 


characterizing key local variations. \textit{IEEE Trans. Image Process.} 


13(6):739--50. 





\bibitem{1-p-1} %2


\Aue{Yu, L., D.~Zhang, and K.~Wang}. 2007. The relative distance of key point 


based iris recognition. \textit{Pattern Recogn.} 40(2):423--30.


\bibitem{2-p-1} %3


\Aue{Hollingsworth, K., K.~Bowyer, and P.~Flynn}. 2009. The best bits in an iris 


code. \textit{IEEE TPAMI} 31(6):964--73.





\bibitem{4-p-1}


\Aue{Pavelyeva, E.\,A., and A.\,S.~Krylov}. 2010. Poisk i analiz kljuchevyh tochek 


raduzhnoj obolochki glaza metodom preobrazovaniya Hermita [Hermite transform 


based iris key points selection and analysis]. 


\textit{Informatics and Applications} 4(1):79--82.


\bibitem{5-p-1}


\Aue{Oppenheim, A.\,V., and J.\,S.~Lim.} 1981.


The importance of phase in signals. 


\textit{IEEE Proceedings} 69(5):529--41.


\bibitem{6-p-1}


\Au{Miyazawa, K., K.~Ito, T.~Aoki, K.~Kobayashi, and H.~Nakajima}. 2006.


A~phase-based iris recognition algorithm. \textit{Advances in biometrics}. 


Eds.\ D.~Zhang and A.\,K.~Jain. Lecture notes in computer science ser.


Berlin, Hridelberg: Springer. 


3832:356--65.


\bibitem{7-p-1}


\Aue{Nagashima, S., K.~Ito, T.~Aoki, H.~Ishii, and K.~Kobayashi}. 2009. 


High-accuracy estimation of image rotation using 1D phase-only correlation. \textit{IEICE 


Trans. Fund. Electr.} E92-A(1):235--43.


\bibitem{8-p-1}


\Aue{Pavelyeva, E.\,A.} 2012. Metod proekcionnoj fazovoj korreljacii v kljuchevyh 


tochkah raduzhnoj obolochki glaza [Hermite projection phase-only correlation 


method in iris key points].  \textit{GraphiCon'2012 Conference 


Proceedings: 22nd Conference (International) on Compuder Graphics and Vision}. 


Moscow. 128--32.


\bibitem{9-p-1}


\Aue{Martens, J.-B.} 1990. The Hermite transform-theory. \textit{IEEE Trans. 


Acoust. Speech} 38(9):1595--606.





\pagebreak





\bibitem{10-p-1}


\Aue{Pavelyeva, E.\,A., and A.\,S.~Krylov}. 2011. Opredelenie lokal'nyh sdvigov 


izobrazhenij raduzhnyh obolochek glaz metodom proekcionnoj fazovoj korreljacii 


[Local iris images shifts detection using Hermite projection phase-only correlation 


method]. \textit{GraphiCon'2011 Conference Proceedings: 


21nd  Conference (International) on Compuder Graphics and Vision}. Moscow. 


188--91.





 \label{end\stat}


 


\bibitem{11-p-1}


\Aue{Krylov, A., and D.~Korchagin}. 2006. Fast Hermite projection method. 


\textit{Image analysis and recognition}. Eds.\ A.~Campilho and M.\,S.~Kamel.


Lecture notes in computer science ser. Berlin, Heidelberg: Springer-Verlag. 4141:329--38.


\bibitem{12-p-1}


CASIA-IrisV3 database. Available at: {\sf http://www.cbsr.ia.ac.cn/IrisDatabase.htm}


(accessed October 25, 2013).


 \end{thebibliography}


} }





\renewcommand{\bibname}{\protect\rmfamily Литература}


